\subsubsection{模型的建立}\label{sec:q3-2}

\textbf{（1）决策内容与目标。}要决定的是机器狗每一步去哪里、在那里测哪些频道、什么时候对哪个源做光学清除，目标是在清除全部源的前提下让式\eqref{eq:q3-2} 的总耗时尽量小。真实源的数量和位置事先不知道，每一次检测的结果又会改变后面该做什么，所以不可能一次把整局的路线规划出来。本文的做法是分层：先用一组固定的检测站保证任何源都会被发现；对已发现的源用定位区域描述它可能在哪里，按第二问的模型选点测向，区域很小时先试一次光学清除；最后把“去哪个站”和“处理哪个源”合成一个路线问题，每走一步重新规划。下面依次建立这四部分。

\textbf{（2）发现覆盖：$7$ 个检测站。}接收半径至少 $1000$ m，所以只要某个位置离站不到 $1000$ m，站在那里扫一遍所有未发现频道，这个位置上的源就一定会被收到。要保证场地里任何位置的源都能被发现，只需一组站，使它们的 $1000$ m 圆盘把整个半径 $1800$ m 的圆域盖住。本文取半径 $999$ m 的圆周上等角分布的 $7$ 个站，见图\ref{fig:q3-cover}(a)。相邻两站相隔 $360^\circ/7\approx51.4^\circ$，圆周上离站最远的点在两站中间，它到最近站的距离约 $999$ m，刚好在 $1000$ m 以内；原点离每个站都是 $999$ m，也在圆内；场地内部的点更近。程序不靠抽样检查覆盖，而是用圆弧关系做一次完整的核验：先检查场地圆周是否全被各站圆盘盖住，再对每个站检查它的圆在场地内的那段弧是否全被其他站的圆盘盖住。如果有一块地方没被盖住，它的边界上一定有某个站圆的一段露出来的弧，所以两条都通过就说明覆盖完整。

这个布局有一点和直觉不同：机器狗从原点出发，却不先在原点把 $20$ 个频道扫一遍。原因是原点在每个站的 $1000$ m 圆内，$7$ 个站扫过之后原点附近的源自然会被发现，而在原点多扫一遍要花 $20\times5+19$ s 左右。开局先扫原点的做法作为对照方案保留在结果分析里，实验表明它平均更慢。

固定的 $7$ 站是计划，不是必须逐一执行的。扫过几站以后，剩下的站可能已经多余：如果去掉某个未去的站，已扫描的站和其余未去的站仍能通过覆盖核验，就把它剪掉；每次只剪一个，剪完再核验。未去的站也可以挪：把它往当前路线的投影点、路线中点或附近已知源的方向试探，只要挪过去以后覆盖核验仍然通过，而且局部路程至少缩短 $1$ m，就采用新位置，见图\ref{fig:q3-cover}(b)。剪站和挪站都只改计划，不改已经扫过的记录，所以“未发现频道判为无源”的依据始终是真实扫描过的位置。

\begin{center}
\begin{minipage}{\linewidth}
\centering
\QThreeCoverFigure
\captionof{figure}{发现覆盖的布局与调整。(a) 半径 $999$ m 圆周上的 $7$ 个检测站，每站的 $1000$ m 接收圆合起来盖住整个场地。(b) 扫过几站之后，负责区域已被其他站盖住的站剪掉，未去的站往路线附近挪，两种调整都要重新通过覆盖核验。}
\label{fig:q3-cover}
\end{minipage}
\end{center}

\textbf{（3）已发现源的定位：选点规则。}某个频道收到第一次示向度后，定位区域是一条长约 $1500$ m 的窄带，处理它时按第二问的思路继续选点测向。本问对第二问的候选点作了调整，以适应多次测向和从任意位置出发的情形。设当前位置为 $x_0$，区域包围圆的圆心为 $z$、半径为 $r$，$u$ 为从 $x_0$ 指向 $z$ 的单位向量，$v$ 与 $u$ 垂直。候选点取圆心 $z$、当前位置 $x_0$，以及
\[
z+r\,(a\,u+b\,v),\qquad a\in\{0,\pm0.35,\pm0.7\},\ b\in\{0,\pm0.15,\pm0.4\}
\]
中除圆心以外的 $24$ 个点，共至多 $26$ 个，见图\ref{fig:q3-probe}(a)。横向偏移 $b$ 的作用和第二问一样，是拉开交会角；纵向偏移 $a$ 让候选点有靠近当前位置、少走路的选择。去掉该频道已经测过的位置，再去掉不能保证收到信号的点，即到区域任一顶点的距离达到 $999.9$ m 的点。

每个候选点 $x$ 的代价仍是“移动时间加测后剩余代价”。测后区域取决于源的真实位置和这次的角误差，两者都未知，处理办法是取平均：源的位置按面积在区域内取样（把多边形剖成三角形，每个三角形取三个内点，按面积加权），角误差取 $-0.77^\circ$、$0$、$+0.77^\circ$ 三个值，按 $5/18$、$4/9$、$5/18$ 加权，这是 $\pm1^\circ$ 区间上三点求积的节点和权重。对每一对假想的源位置 $p$ 和角误差，把区域与新楔形求交，得到新的包围圆圆心 $z'$ 和半径 $r'$，剩余代价按下式估计：
\begin{equation}\label{eq:q3-4}
P(x)=
\begin{cases}
\dfrac{\max\{0,\ \|z'-x\|-(20-r')\}}{v}+5, & r'\le19.5\ \text{m},\\[8pt]
\dfrac{\|z'-x\|+\max\{0,\ \|p-z'\|-19.5\}}{v}
+6\log_2\dfrac{r'}{19.5}+\dfrac{0.5\,r'}{v}+5, & r'>19.5\ \text{m}.
\end{cases}
\end{equation}
上一行是测完就能清除的情形：走到清除位置再清除。下一行是还要继续测的情形：走到新区域中心、再走到源的时间，加上后续测向的时间，按“每测一次半径大约减半、每次检测加换频 $6$ s”估计为 $6\log_2(r'/19.5)$，再加一项与半径成比例的额外移动。代价取所有假想情形的加权平均，再加上从 $x_0$ 到 $x$ 的移动时间和 $5$ s 检测，选代价最小的候选点。

这一套选点每个源最多做 $4$ 次；超过 $4$ 次，或者包围圆半径已经小于 $22$ m，就直接走到包围圆圆心测一次。在圆心测向有确定的效果：源在圆心的某一侧，楔形从圆心出发，新区域落在一条从圆心出发、长不超过 $r$ 的窄条里，包围圆半径至少减半。所以无论前面的候选点选得好不好，一个源总能在有限次测向内达到清除条件。圆心离区域里每个点都不超过 $r$，远小于 $1000$ m，在那里一定能收到信号。

\textbf{（4）清除判据与清除位置。}光学清除要求离源不超过 $20$ m。定位区域的包围圆半径 $r\le r_c=19.5$ m 时，所有到区域每个顶点都不超过 $20$ m 的位置构成一个非空的凸集，圆心就在其中；走到其中任何一点清除都一定成功。程序在这个集合里选一个点，使“从当前位置走过去”加上“从那里走到下一个任务”的总路程最短：候选包括从当前位置和从下一任务位置各自最近的可行点、两者连线上的插值点，以及当前位置到下一任务的线段与可行集的交（若相交，清除就不需要绕路）。此外，每一步开始时先检查当前位置：若它离某个已知源区域的所有顶点都不超过 $20$ m，就地清除，不再移动。

\textbf{（5）提前光学清除。}上面的清除判据是充分条件，要求区域缩到半径 $19.5$ m 以内，这往往需要在很近的距离上再测一两次。一次失败的光学尝试只花 $3$ s，比一次检测加换频的 $6$ s 还便宜，而且机器狗常常已经站在区域附近。于是本文加了一条规则：区域还没有小到能保证清除，但一个 $20$ m 圆盘已经能盖住它的大部分时，先试一次光学清除，失败了再按原来的流程测。

“盖住大部分”用面积占比来衡量。把定位区域剖成三角形，每个三角形再分成 $10\times10$ 个等面积的小三角形，取各自的重心作为面积样本点 $x_j$，权 $w_j$ 与所在三角形的面积成正比、总和为 $1$。位置 $q$ 处 $20$ m 圆盘盖住的面积占比为
\begin{equation}\label{eq:q3-14}
h_f(q)=\sum_j w_j\,\mathbf 1\{\|x_j-q\|\le20\},
\end{equation}
见图\ref{fig:q3-probe}(b)。它只是区域面积的比例，不是源落在圆盘内的概率，因为源在区域内的分布并不知道；用它来决定要不要花 $3$ s 赌一次，是够用的。提前尝试的条件有四条：该频道正在定位中，且包围圆半径不超过 $80$ m；对这个频道的提前尝试还不到 $2$ 次，成功和失败都计数；当前坐标没有对它试过；$h_f(\text{当前位置})\ge0.4$。四条都满足才向设备发出清除请求。请求成功，该频道转入已清除；失败，定位区域和所有状态原样保留，不做任何推断，因为失败的原因可能是源在圆盘外的那部分区域里，而这部分区域本来就没有被排除。

提前尝试还改变选点：只要一个源满足前三条资格、局部测向还不到 $4$ 次，就不再按（3）的代价选点，而是走到区域的面积加权重心 $q^*=\sum_j w_jx_j$，条件是 $h_f(q^*)\ge0.4$、那里能保证收到信号、且没有对它试过。到达以后先试一次清除；失败再在原地测向、更新区域，然后继续。每源最多两次尝试，所以失败的光学费用每源不超过 $6$ s；但去重心和去代价最小点走的路不同，总的影响要靠实验看。

\begin{center}
\begin{minipage}{\linewidth}
\centering
\QThreeProbeFigure
\captionof{figure}{定位阶段的两种选点。(a) 在包围圆圆心附近按纵向、横向偏移布设候选探测点，超出可靠接收范围或已测过的点剔除，按式\eqref{eq:q3-4} 的代价取最小。(b) 区域已经较小时，在面积加权重心处一个 $20$ m 圆盘盖住的面积占比 $h_f$ 达到 $0.4$，就先试一次光学清除。}
\label{fig:q3-probe}
\end{minipage}
\end{center}

\textbf{（6）任务的安排与共享测量。}每一步的任务集由两部分组成：还没去的检测站，位置就是站的坐标；已发现但没清除的源，位置用当前包围圆的圆心代表。把当前位置当起点，对这些任务求一条不必返回的最短路线，只执行路线上的第一个任务，并记下第二个任务的位置供（4）选清除点用；执行完一个任务，区域和任务集都变了，再重新排。任务最多 $7+16$ 个，程序用多起点的局部搜索求路线：以每个任务为第一站做最近邻排序，各自用 2-opt 和单点移位改进到不能再改，再对最好的一条随机交换两三对任务后重新改进，重复 $40$ 次，取最短的。每次规划只要几毫秒。

机器狗到一个站扫完未发现频道之后，人已经在那里，对附近的已知源顺便再测一次往往划算。程序对每个正在定位的源做一次估计：假设从当前位置对它测向、示向度正指区域中心，算出新区域的包围圆半径；若能缩到原来的 $0.8$ 倍以下并且至少缩小 $30$ m，或者直接缩到清除门限以内，就补测一次，前提是当前位置离这个源不超过 $1500$ m 加半径，且没有在这里测过。清除一个源以后、对某个源测向以后，也用同样的规则检查其他源。另外，如果在当前位置把未发现频道扫一遍能让某个计划中的站被剪掉，就在这里扫，省去以后专门跑一趟。

\textbf{（7）终止条件。}任务在两种情况下结束：已经清除了 $16$ 个不同频道，源数上限保证没有遗漏；或者计划中的站全部扫完（含被剪掉的站，剪站时已核验其余站足以覆盖），此时每个未发现频道在场地任何位置都会被某个已扫描的站收到，可以判为无源，再把已发现的源全部清除。程序在两种情况下都检查 $20$ 个频道是否全部落在已清除或无源两类，缺一个都不退出。发现的源已达 $16$ 个但还有源没清除，或者站扫完了还有已知源，都不能结束。
